%% Le lingue utilizzate, che verranno passate come opzioni al pacchetto babel. Come sempre, l'ultima indicata sarà quella primaria.
\def\thudbabelopt{italian,english} %% "english" is placed last to make it the primary language for the text.
%% Valori ammessi per target: bach (tesi triennale), mst (tesi magistrale), phd (tesi di dottorato).
%% Valori ammessi per aauheader: aics (Department of Artificial Intelligence and Cybersecurity)
\documentclass[target=mst,aauheader=aics,style=]{thud}

%% --- Informazioni sulla tesi ---
\course{Cybersecurity and Artificial Intelligence}
\title{Automated Threat Modeling Pipeline for Operational Technology Configurations using Declarative Logic Synthesis}
\author{Xhulia Palaj}
\supervisor{Prof.\ Marino Miculan}
\cosupervisor{Dr.\ Jasmin Wachter}
\tutor{Dott.\ Canzian}

%% --- Pacchetti consigliati ---
\usepackage[a-1b]{pdfx}
\usepackage[pdfa]{hyperref}
\usepackage{amsmath}
\usepackage{booktabs}

\begin{document}
\maketitle

%% Dedica (opzionale)
\begin{dedication}
    To my family and friends.
\end{dedication}

%% Ringraziamenti (opzionali)
\acknowledgements
I would like to express my deepest gratitude to my supervisor, Professor Marino Miculan, and my co-supervisor, Dr. Jasmin Wachter, for their invaluable guidance, constant support, and immense patience throughout the development of this research project. I am equally grateful to my industrial tutor, Dott. Canzian, for providing real-world operational insights and configuration baselines that made this study practically viable. Finally, a special thank you goes to my peers at the student dormitory for their encouragement and support during these months.

%% Sommario (opzionale)
\abstract
Modern industrial automation architectures rely heavily on the convergence of Information Technology (IT) and Operational Technology (OT). While this integration optimizes industrial processes, it introduces severe cyber risks by exposing legacy, unauthenticated communication protocols to external threat landscapes. Standard active vulnerability evaluation tools are unusable in live OT safety environments due to the risk of triggering unpredictable kinetic failures or operational downtime. 

To bridge this validation gap, this thesis research presents an automated, non-intrusive threat modeling pipeline designed to parse raw semi-structured infrastructure asset parameters directly into mathematically sound Logical Attack Graphs (LAGs). Leveraging custom-designed Datalog Horn clauses within a monotonic logic framework based on MulVAL, our pipeline integrates the strict architectural constraints of the Purdue Reference Model and maps structural access vulnerabilities onto the MITRE ATT&CK for ICS framework. Drawing structural inspiration from the micro-level firmware execution abstractions pioneered by the MIRAGE project, this model scales logic-based verification upward to evaluate multi-stage lateral movement risks across enterprise network boundaries. To bypass the historical bottlenecks of manual asset translation, we evaluate advanced synthesis paradigms, utilizing Large Language Models (LLMs) verified via automated, multi-tiered linting and deterministic validation layers. The framework guarantees robust risk prioritization with zero production disruption.

\tableofcontents

%% Corpo principale del documento
\mainmatter

% ==========================================
% CHAPTER 1: INTRODUCTION
% ==========================================
\chapter{Introduction}

\section{The Industrial Paradigm Shift \& Contextualizing OT Constraints}
Industrial Control Systems (ICS) and Operational Technology (OT) networks form the critical infrastructure backbone of manufacturing plants, power grids, and heavy industrial operations. Historically isolated through physical separation or "air-gapping," modern industrial competitive demands have forced a widespread convergence between Information Technology (IT) enterprise networks and OT production environments. However, directly applying IT security frameworks to OT ecosystems introduces profound systemic friction due to fundamentally opposing operational priorities.

In traditional IT environments, cybersecurity paradigms are governed by the classic CIA Triad (Confidentiality, Integrity, and Availability), prioritizing the protection of data data-at-rest and data-in-transit above all else. Within Operational Technology, this hierarchy undergoes a complete inversion, transitioning into a paradigm focused on continuous availability, predictability, and human/physical safety.

An unauthorized data leak in an IT environment represents a severe financial or regulatory violation, but a forced system shutdown in an OT environment can trigger catastrophic kinetic failures, irreversible equipment damage, and multi-million-dollar operational losses. Consequently, cybersecurity strategies within industrial sectors cannot rely on reactive patch management or intrusive security operations that threaten systemic uptime. 

\section{The Economic and Physical Risks of Active Vulnerability Testing}
Because industrial assets, such as Human-Machine Interfaces (HMIs), Programmable Logic Controllers (PLCs), and Supervisory Control and Data Acquisition (SCADA) servers, frequently run proprietary, legacy software protocols, they exhibit acute structural fragility. Standard IT risk management relies heavily on active network scanners (e.g., Nessus, Nmap) to ping ports and query software banners for active vulnerability mapping. In an OT network, the high packet traffic and unexpected request formats generated by these active scanners can overwhelm the fragile memory buffers of legacy embedded field devices, forcing PLCs into fatal fault modes or causing temporary control-loop blindness.

Given that a single minute of unexpected downtime in heavy industries like steel manufacturing, automotive assembly, or chemical processing can incur enormous financial penalties, live penetration testing, and active vulnerability scanning cannot be performed arbitrarily. Instead, comprehensive risk assessments, rigorous modeling, and formal test verification plans must be completed completely out-of-band before an industrial network configuration is commissioned or updated in a live plant environment.

\section{Limitations of Manual Risk Analysis in Large-Scale Industrial Automation}
Modern heavy industrial automation architectures, such as those deployed across complex manufacturing facilities, are massive in scale, containing highly segmented networks that mirror the traditional Purdue Reference Model. A single plant deployment frequently integrates:
\begin{itemize}
    \item Dozens of virtualized environments and virtual machines handling process historical records, engineering updates, and localized server tasks.
    \item Distributed Human-Machine Interfaces (HMIs) interacting across diverse physical zones.
    \item Redundant Active Directory Domain Controllers managing plant-wide access configurations.
\end{itemize}

Conducting manual risk assessments over an infrastructure of this scale is mathematically and logistically inefficient, making it highly prone to human error. Manual evaluation teams often analyze assets in isolation, failing to account for the complex interdependencies created by cross-zone firewall policies, system configurations, and shared administrative user accounts. 

To overcome these manual limitations, the automated generation of Logical Attack Graphs (LAGs) is essential. By building automated assessment pipelines, security architects can formally mathematically model and prove how a remote adversary might navigate through network perimeters, discover configuration flaws, exploit legacy protocol vulnerabilities, and achieve lateral movement to compromise critical physical processes.

\section{Objectives and Core Scope of the Thesis Research}
The primary objective of this research is to evaluate the real-world effectiveness of an automated threat modeling pipeline designed to ingest raw industrial infrastructure metrics, translate them into formal logic variables, and generate predictive attack paths.

This study aims to establish a multi-phased pipeline that bridges the gap between human-readable data structures and rigorous mathematical validation:
\begin{enumerate}
    \item \textbf{Asset and Configuration Ingestion:} Extracting host parameters, software service descriptions, and operating system properties directly from machine-readable specifications.
    \item \textbf{Intermediate Data Structuring:} Normalizing and parsing raw, disparate asset metadata into a clean, unified JSON serialization schema.
    \item \textbf{Formal Logic Translation:} Evaluating advanced parsing approaches, including Large Language Models (LLMs), to automate the translation of structured network states into formal declarative Datalog facts ready for ingestion by the MulVAL reasoning engine.
    \item \textbf{Threat Intelligence and Hardening Integration:} Mapping generated logical attack sequences directly onto real-world adversary tactics and techniques using the MITRE ATT&CK for ICS framework to output clear, actionable hardening guidelines.
\end{enumerate}

Furthermore, this study looks directly into the underlying mechanics of automated logic frameworks, exploring whether automated parsing models introduce systemic biases or exhibit structural visibility gaps when evaluating specialized industrial attack vectors compared to traditional IT vulnerabilities. Through empirical testing, this thesis aims to deliver a scalable, production-grade model to verify complex OT infrastructures securely, efficiently, and with zero operational disruption.


% ==========================================
% CHAPTER 2: BACKGROUND
% ==========================================
\chapter{Background}

\section{Explanation on Attack Trees and Graphs \& Logic Representations}

\subsection{The Limitations of Traditional Graph-Theoretic State-Space Models}
The early formalization of automated security analysis in networked systems relied heavily on traditional state-space transition systems. In these models, cybersecurity risk assessment was treated as a verification problem where the entire network environment was modeled as a finite state machine. Each node in the graph represented an explicit, global state of the network, comprising the exact configuration, patch levels, operational parameters, and privileges across every single host simultaneously. Edges between these nodes represented discrete atomic transitions, such as an adversary executing a specific exploit to move from one state to another.

While conceptually intuitive, state-space models suffer from a catastrophic mathematical limitation known as state-space explosion. Because these frameworks attempt to track every possible permutation of a network's global state, the size of the graph scales exponentially with the number of hosts, services, and vulnerabilities. For a network with $H$ hosts, each capable of being in $S$ distinct security configurations, the potential state space expands to $O(S^H)$.

When applied to large-scale infrastructure environments characterized by hundreds of virtual machines, distributed Human-Machine Interfaces (HMIs), and complex enterprise-wide asset footprints, computing the full state transition graph quickly exceeds reasonable processing constraints. This computational bottleneck makes real-time, dynamic risk assessments mathematically unfeasible under a state-space architecture.

\subsection{Shift to the Monotonic Logic Paradigm and Horn-Clause Datalog}
To bypass the exponential overhead of state-space exploration, modern predictive security modeling transitioned to a formal logic paradigm. Pioneered by frameworks like MulVAL (Multi-host, Multi-stage Vulnerability Analysis Language), this approach changes how threat propagation is computed. Instead of tracking explicit global system states, a logic-based model views the network through a declarative lens, evaluating the environment using mathematical clauses written in a subset of first-order logic known as Horn-clause Datalog.

This paradigm shift relies fundamentally on the principle of monotonicity. In a monotonic logic framework, it is assumed that an adversary's gained privileges accumulate over time and are never lost during an ongoing attack scenario. If an attacker gains administrative control over a specific server, that privilege remains a true fact within the engine's logical environment throughout the duration of the analysis path. By treating security analysis as a monotonic derivation problem, the logic engine does not need to compute or track a volatile, time-varying graph of explicit system configurations. Instead, the complexity is drastically reduced to polynomial time $O(N^2)$, where $N$ represents the number of independent configuration facts and security parameters in the environment.

\subsection{Taxonomies of Automated Attack Generation Methodologies}
In their comprehensive methodological surveys, Konsta et al. establish a clear taxonomic classification that organizes predictive security modeling frameworks into three primary generation paradigms:
\begin{itemize}
    \item \textbf{Pure Graph-Theoretic Architectures:} These systems construct explicit directed graphs to map out paths, relying on classic traversal algorithms to evaluate reachability. While highly effective for small, well-defined environments, they are fundamentally limited by the state-space explosion problem when scaled horizontally across enterprise topologies.
    \item \textbf{Model-Checking Frameworks:} Leveraging verification systems such as NuSMV or SPIN, these approaches model network components as strict state machines and use temporal logic formulas (e.g., LTL or CTL) to verify security properties. While mathematically rigorous and highly capable of tracking dynamic system states, they require substantial processing power and often face scalability issues in large infrastructures.
    \item \textbf{Rule-Based Declarative Logic Engines:} Frameworks such as MulVAL leverage Horn-clause logic programming and backward-chaining engines to evaluate security risks. This architecture strikes an effective balance for large enterprise and industrial footprints: it achieves high scalability and expressiveness by abstracting network states into static relational facts and generic interaction rules. However, its core limitation lies in its static nature; because it assumes a monotonic environment, it natively struggles to model temporary states, complex temporal timing conditions, or active cyber-physical feedback loops without specialized domain extensions.
\end{itemize}

\subsection{Algorithmic Path Discovery \& Optimization Heuristics}
As industrial automation environments scale horizontally, computing the complete logical derivation of an attack graph becomes a significant computational challenge. While formal logic programming engines like MulVAL successfully avoid the exponential explosion found in traditional state-space models by executing polynomial-time backward-chaining algorithms $O(N^2)$, constructing a comprehensive Logical Attack Graph (LAG) across an enterprise footprint can still introduce major processing overhead. 

When a network incorporates hundreds of segmented virtual subnets, intricate zone-to-zone firewall parameters, and widespread asset dependencies, compiling the entire permutation of every theoretically possible attack path generates massive rule dependency matrices. Every additional node or credential relationship expands the logical derivation tree, increasing execution times and memory utilization during full database compilations. This processing bottleneck becomes particularly acute when conducting real-time risk assessments, where security architectures must adapt dynamically to rapid topology alterations or newly discovered software vulnerabilities. To maintain responsiveness in large-scale industrial infrastructures, contemporary security engineering shifts from exhaustive graph compilation toward optimized path discovery techniques.

\subsubsection{Heuristic Pathfinding Models vs. Reinforcement Learning}
To mitigate compilation overhead, modern security frameworks focus on optimizing path discovery, primarily utilizing two distinct approaches: Machine Learning (specifically Deep Reinforcement Learning) and Cost-Aware Heuristic Search Algorithms. 

\paragraph{Deep Reinforcement Learning (DRL)}
DRL models approach path optimization by deploying autonomous agents that learn optimal adversarial propagation pathways through continuous environment interaction and reward maximization. While highly capable of identifying complex attack paths across volatile topologies, DRL models demand substantial offline training phases, massive datasets, and intensive hardware resources. Furthermore, the "black-box" nature of deep neural networks introduces uncertainty, making it difficult to extract clear, deterministic mathematical explanations for the derived paths, a major drawback for verification procedures in safety-critical industrial settings. 

\paragraph{Cost-Aware Heuristic Search Variants}
To sidestep the resource-heavy training demands of machine learning, modern research implements real-time, deterministic path discovery frameworks like the Automated Attack Path Discovery and Exploitation Model (AAPDEM). Rather than compiling the complete attack graph or running intensive neural training cycles, AAPDEM utilizes an optimized, cost-aware $A^*$ search variant to isolate high-risk path configurations incrementally. 

The algorithmic execution defines an evaluation function:
\begin{equation}
f(n) = g(n) + h(n)
\end{equation}
where $g(n)$ represents the accumulated computational cost or exploitation difficulty an adversary faces when traveling from the initial perimeter entry point to the current logical node $n$. The parameter $h(n)$ represents a directional heuristic weight that estimates the structural distance, network hops, and exploitation barriers remaining between node $n$ and the ultimate industrial target (e.g., a Level 1 PLC). By using this heuristic function to prune low-probability or high-difficulty execution paths early, the framework can process network changes in real time, dramatically lowering processing demands while maintaining accuracy. 

\subsubsection{Graph Neural Networks (GNNs) for Temporal and Formal Security Properties}
While cost-aware heuristics like the $A^*$ algorithm optimize static path discovery, evaluating complex temporal behaviors and verifying formal safety properties across multi-stage graphs has prompted exploration into Graph Neural Networks (GNNs). Traditional logic frameworks evaluate environments through monotonic assumptions, meaning they are natively poorly suited for tracking transient environmental states, time-dependent firewall configurations, or reactive cyber-physical feedback loops. 

To address these limitations, advanced research tracks, such as work analyzing Büchi Automata with Graph Neural Networks, leverage neural structures specifically designed to process graph-structured data. GNN architectures use continuous message-passing mechanics to aggregate relational features from neighboring nodes directly, allowing them to capture non-linear node dependencies across complex topologies. 

When applied to formal verification frameworks, a trained GNN can analyze underlying safety automata or logical attack trees to predict whether a specific multi-stage attack sequence violates a plant's temporal safety rules. This development paves the way for hybrid threat modeling tools: formal verification engines handle the structured compilation of static asset facts and interaction rules, while intelligent graph neural network layers accelerate or optimize path discovery to expose hidden, time-dependent security risks.

\section{MulVAL}
The core execution architecture of MulVAL achieves high scalability by separating the threat model into two distinct, decoupled components: Static Configuration Facts and Generic Interaction Rules.

\subsection{Static Configuration Facts}
These represent the immutable, factual reality of the target network environment. They are expressed as Datalog clauses that register network topology, firewall access control lists (ACLs), software installations, and software vulnerabilities. Typical predicates include:
\begin{itemize}
    \item \texttt{located(Host, Subnet, Type)}: Maps the network topology location of an asset.
    \item \texttt{networkService(Host, Program, Protocol, Port, Privilege)}: Defines a service binding and its local run permissions.
    \item \texttt{vulHost(Host, CVE, Program, ExploitType, Consequence)}: Maps known vulnerability intelligence onto specific hosts.
    \item \texttt{aclNW(Source, Destination, Protocol, Port)}: Represents network filtering and firewall rules.
\end{itemize}

\subsection{Generic Interaction Rules}
These rules act as the underlying mathematical "adversarial physics" that dictate how an attacker can navigate through an environment. These are written as generic, abstracted Horn clauses that define under what precise conditions a new logical deduction can be derived from existing facts. For example, a standard remote exploit interaction rule is structured as:
\begin{equation*}
\text{execCode(Target, Perm)} \leftarrow \text{networkService(Target, Prog, Prot, Port, Perm)}, \dots
\end{equation*}

\subsection{The Derivation Engine and Table Optimization}
MulVAL utilizes an underlying XSB Prolog execution engine to perform a goal-directed, backward-chaining derivation. Given a user-defined target condition, such as \texttt{domainCompromise('Danieli\_Domain')}, the engine works backward, evaluating which combinations of interaction rules and configuration facts can satisfy the terminal objective.

To optimize this process and prevent infinite loops caused by cyclic network dependencies (such as two hosts running identical software and sharing reciprocal network access), the engine leverages XSB Tabling Optimization Directives (e.g., \texttt{:- table execCode/2.}). Tabling enforces a memorization strategy: once a specific security subgoal is evaluated, its truth value is stored in an internal lookup table. If the engine encounters the same subgoal later in another branch of the attack tree, it pulls the cached result instantly rather than re-evaluating the entire logic sequence. This ensures that the engine can dynamically construct a comprehensive, multi-stage Logical Attack Graph (LAG) while keeping computation times low.

\section{MIRAGE}

\subsection{Multi-Binary Firmware Vulnerability Context}
While traditional logic-based risk assessment frameworks focus on computing vulnerabilities at the macroscopic enterprise network or physical topology layers, a novel boundary in automated security analysis shifts attention to a device’s internal, microscopic ecosystem. Modern Internet of Things (IoT) and Operational Technology (OT) edge devices heavily utilize lightweight, multi-binary operating structures. In these architectures, complex operational tasks—ranging from web-based management interfaces to external industrial protocol parsing—are distributed across disparate system binaries that communicate locally via Inter-Process Communication (IPC) channels. 

As established by Tayouri et al., evaluating vulnerabilities in complete isolation fails to mirror realistic multi-stage adversarial pathways within device firmware. A single vulnerability within a non-network-facing internal binary can have systemic security impacts if an adversary can execute an intranode pivot. By exploiting internal data flows, shared environment variables, or execution chains, an attacker can propagate laterally between system binaries, effectively escalating privileges or hijacking the entire device from an otherwise secure boundary. 

\subsection{The MIRAGE Framework Architecture}
To bridge the gap between isolated vulnerability scanning and realistic internal threat vectors, Tayouri et al. introduced the MIRAGE (Multi-Binary Image Risk Assessment with Attack Graph Employment) framework. MIRAGE pioneers the application of Logical Attack Graphs (LAGs) explicitly to multi-binary firmware ecosystems by implementing a static analysis pipeline operating across four sequential stages: 
\begin{enumerate}
    \item \textbf{Binary Modeler:} This module decompresses compressed firmware filesystems using an abstract tool layer (defaulting to Binwalk) to chart border (network-facing) binaries and map inter-binary data dependencies. 
    \item \textbf{Vulnerability Mapper:} This component filters the NIST National Vulnerability Database (NVD) using static indicator heuristics, token string matches, and version-querying command execution under QEMU emulation to generate precise \texttt{vulExists} clauses for MulVAL ingestion. 
    \item \textbf{Attack Graph Generator:} This stage extends MulVAL’s core reasoning engine with specialized firmware-related interaction rules tailored to trace external network boundaries and evaluate internal supply chain infections via open-source software (OSS) bugs. 
    \item \textbf{Attack Graph Analyzer:} Incorporating both a Binaries Risk Analyzer and an Attack Paths Risk Analyzer, this module enumerates attack sequences using a modified successive path algorithm, applying Exploitation Prediction Scoring System (EPSS) metadata to calculate granular risk scores. 
\end{enumerate}

\subsection{Practical Evaluation Insights and Performance Baselines}
The experimental validation of the MIRAGE framework over a robust empirical dataset of 1,343 multi-binary firmware images yielded critical baseline conclusions for automated threat modeling: 
\begin{itemize}
    \item \textbf{High Attack Path Prevalence:} Out of the evaluated dataset, 84\% of firmware images generated full logic attack graphs, verifying that standard firmware builds routinely contain interconnected vulnerable execution boundaries. 
    \item \textbf{Discovery Scale Over Prior Art:} MIRAGE demonstrated a significant expansion in vulnerability visibility compared to older static interaction tracking systems like Karonte. While Karonte restricts its path detection to single-hop boundary loops, MIRAGE identified a dataset average of 14 complex attack paths per image under pure external threats, and 27.4 paths when integrating internal threat hypotheses. 
    \item \textbf{Drastic Performance Optimization:} The logic-based abstraction engine implemented by MIRAGE achieved a 57\% average time performance reduction across wide database scans compared to Karonte's traditional taint-tracking engine, completing parallel pipeline execution tasks in under 2.5 hours on average. 
    \item \textbf{Risk Prioritization Effectiveness:} Through the application of binary interaction counting to measure impact, combined with EPSS probability matching to determine likelihood, the framework successfully highlighted the riskiest foundational components across common distributions (such as identifying the \texttt{file} and \texttt{pg} utilities as major internal propagation nodes). 
\end{itemize}

\section{OT Networks and the Purdue Model (Peculiar Characteristics)}
To securely segment fragile automation equipment from volatile business enterprise networks, industrial networks are strictly structured according to the Purdue Reference Model for ICS Security (ANSI/ISA-95). This hierarchical model organizes the industrial plant infrastructure into distinct, segmented operational layers.

Traditional logic engines treat network topologies as flat, interconnected graphs or simple, direct source-to-destination paths. To build a mathematically precise model of an industrial plant, the Datalog engine must formally enforce these Purdue boundaries using explicit structural facts. Within an operational network environment, this layout is declared by mapping distinct subnets to specific layers using topological predicates:
\begin{verbatim}
located(it_subnet, 'IT_Network', ipSubnet).
located(ot_dmz_subnet, 'OT_DMZ', ipSubnet).
located(l2_subnet, 'Level_2_Network', ipSubnet).
located(l1_subnet, 'Level_1_Network', ipSubnet).
\end{verbatim}

By explicitly declaring these network boundaries, the extended reasoning engine can track perimeter-crossing rules. This allows it to evaluate how an external adversary positioning themselves on a compromised boundary asset inside the DMZ can pivot across network filters down into Level 2 supervisory zones or Level 1 control environments. 

\section{Domain-Specific Logic Extensions for CP OT Systems}
While foundational logic-based vulnerability analysis engines like MulVAL revolutionized automated threat modeling, their native, out-of-the-box rule sets are inherently constrained to traditional IT abstractions. In a standard IT security model, system interactions are generalized around standard operating system permissions (e.g., user, admin, root), network connections established over standard TCP/IP sockets, and software vulnerabilities cataloged almost exclusively via the Common Vulnerabilities and Exposures (CVE) index.

However, this abstraction framework fails when confronted with the unique operational and architectural realities of Operational Technology (OT) and Cyber-Physical Systems (CPS). Industrial environments feature thousands of embedded field devices that lack conventional operating systems or user access control management entirely. In these networks, a compromise does not merely manifest as a privilege escalation from a standard user to a system administrator; instead, it frequently involves the unauthorized injection of unauthenticated control commands directly into a physical process. To map realistic multi-stage adversarial pathways within industrial infrastructures, formal logic threat models must be extended with domain-specific rules that natively capture network hierarchies, industrial protocol behaviors, and cyber-physical dependencies.

\section{MulVAL's Limitations Regarding OT}
Standard scanners and IT logic engines are completely blind to structural configuration defects and legacy protocol flaws because they assess assets individually and assume that the absence of an unpatched software CVE equates to a secure state. Traditional logic engines treat network configurations as flat spaces, neglecting the cascading consequences that network filtering policies impose on targeted embedded field layers. Furthermore, traditional IT logic structures lack predicates to express the unauthenticated and open landscape of device communications that are pervasive at lower architectural layers.

\section{Possible Extensions and Modifications of MulVAL for OT}
A critical security risk in OT environments stems from the pervasive reliance on legacy automation protocols—such as Modbus/TCP, Profinet, and EtherNet/IP—which were designed decades ago with a disregard for active cryptographic security. These protocols inherently lack link-level authentication, integrity validation, or encryption. Consequently, any network node that achieves physical or logical proximity to an industrial bus can execute entirely exploitless compromises, such as spoofing control messages, rewriting register parameters, or issuing force-shutdown commands to field devices.

To capture these protocol anomalies within a formal logic framework, the threat model must introduce custom interaction rules that represent vulnerabilities in the communication protocol itself, rather than software code errors (CVEs). For instance, a lack of protocol authentication on a Modbus or Profinet link is represented as a static protocol flaw: 
\begin{verbatim}
vulLinkProtocol(sub_process_area_1, vul_modbus_no_auth, 
                modbus, adjacent, impersonateDst). 
\end{verbatim}

When an adversary achieves local access to this unauthenticated network layer, they do not need to exploit a traditional software buffer overflow to gain control. Instead, the logic engine utilizes a domain-specific interaction rule to deduce that an attacker can spoof communication links and impersonate critical control masters:
\begin{verbatim}
interaction_rule(
    (spoofLinkHost(Principal, ImpersonateMaster, FooledSlave, 
                   AttackerHost, deception):-
        vulLinkProtocol(BusID, _VulId, Protocol, adjacent, impersonateDst),
        l2Connection(AttackerHost, FooledSlave, BusID, Protocol, bus),
        isMaster(ImpersonateMaster, BusID),
        malicious(Principal)),
    rule_desc('Spoofing as a bus Master due to lack of auth', 1.0)
).
\end{verbatim}
This extended rule allows the engine to mathematically prove how a compromised engineering workstation can bypass traditional security barriers and manipulate physical field equipment directly over unauthenticated operational buses.

\section{Access Control Matrices and Zone-to-Zone Cross-Talk Logic}
Because active network segmentation via industrial firewalls represents the primary line of defense in an OT infrastructure, the logic model must enforce a strict network filtering matrix. In an industrial plant, network access is highly constrained; standard assets inside a Level 2 supervisory network are restricted from interacting directly with Level 0/1 field equipment, except through explicitly authorized, tightly restricted exceptions.

Extended Datalog models represent this structural access control matrix using explicit network access tokens (\texttt{aclNW/4}), which define the exact paths allowed through firewalls:
\begin{verbatim}
aclNW(hmi_client, 'Sub_Process_Area_1', tcp, 4840).
aclNW(jump_server, plc_l1_1, tcp, 502).
\end{verbatim}

The reasoning engine processes these firewall permissions alongside custom crosstalk rules to analyze lateral movement vectors. If an unauthorized path or a misconfigured firewall rule allows a high-tier asset to communicate with a low-tier control device over an industrial port (such as Modbus TCP port 502), the logic engine evaluates this as an immediate threat vector.

By combining access privileges with protocol flaws, the engine can deduce when an adversary is in a position to execute unauthorized industrial command injections:
\begin{verbatim}
interaction_rule(
    (unauthorizedControl(TargetPLC):-
        execCode(AttackerWorkstation, admin),
        netAccess(_Principal, AttackerWorkstation, TargetPLC, tcp, 502)),
    rule_desc('Unauthorized command injection via unauth protocol', 0.8)
).
\end{verbatim}
Through this rule-based approach, the extended logic engine successfully maps how minor configuration oversights and unauthorized crosstalk pathways combine to create critical threat vectors, providing a rigorous framework to assess and harden industrial networks against multi-stage cyberattacks.


% ==========================================
% CHAPTER 3: ANALYSIS
% ==========================================
\chapter{Analysis}

\section{What We Want to Achieve: Use Case}
The ultimate functional goal of our framework is to automate out-of-band threat modeling for large-scale, segmented industrial architectures without risking the stability of fragile live operational systems. We aim to construct an end-to-end translation and evaluation framework that takes raw, human-managed system configuration files from enterprise repositories and accurately derives complex, multi-stage attack vectors that threaten Level 1 field equipment. 

By scaling logic-based verification models upward from isolated nodes to comprehensive plant topologies, the framework identifies structural flaws—such as unauthenticated legacy protocols, credential reuse across distinct zones, and firewall crosstalk errors. These derived security gaps are immediately mapped to the real-world adversary tactics defined within the MITRE ATT&CK for ICS matrix, providing plant operators with an automated, zero-disruption optimization pipeline to implement proactive hardening strategies.

\section{Requisite: Input/Output of What We Want to Implement}
Developing an automated logic synthesis engine requires overcoming a fundamental architectural mismatch between how semi-structured formats (e.g., JSON or XML) and declarative logic programming languages represent data.

\subsection{System Input (Semi-Structured JSON Configuration Data)}
Consider the structure of a standard industrial asset inventory entry managed within an engineering database, matching the schema provided by the production configuration profile:
\begin{verbatim}
{
    "id": "vJumpServer",
    "os": "Windows Server 2022",
    "credentials": {
        "admin": "Administrator/Automation%1969"
    },
    "software_services":[
        {
            "name":"active_directory", 
            "port":389, 
            "protocol":"tcp"
        }
    ]
}
\end{verbatim}
This semi-structured layout relies on hierarchical encapsulation and object nesting to represent system relationships. While convenient for object-oriented software engineering, this format is completely incompatible with a Horn-clause Datalog reasoning engine.

\subsection{System Output (Relational Logical Predicates)}
To evaluate this asset, a logical derivation engine like MulVAL requires a flat, relational data structure based on mathematical predicates. The hierarchical JSON blocks must be completely broken down, unnested, and synthesized into distinct relational clauses that isolate independent system variables:
\begin{verbatim}
hasAccount(vJumpServer, admin, 'Administrator/Automation%1969').
networkService(vJumpServer, active_directory, tcp, 389, admin).
\end{verbatim}

A deterministic parser built explicitly for a single schema can achieve this translation through rigid string mapping. However, these static parsers lack flexibility; they are brittle and fail when encountering unindexed object structures, schema alterations, or heterogeneous data formats. Consequently, modern research focuses on deploying advanced parsing techniques capable of handling semantic and syntactic variations intelligently.


% ==========================================
% CHAPTER 4: PROJECT
% ==========================================
\chapter{Project}

\section{Inspired from MIRAGE}
The architectural layout and practical conclusions extracted from the MIRAGE paradigm provide a direct roadmap for the macro-environmental automation goals established in this thesis. While MIRAGE operates at a micro-level to map how an attacker propagates between binary execution contexts inside a single device, this research scales that exact logical methodology upward to an internode network infrastructure. 

By migrating the concepts of inter-binary data flow tracking to zone-to-zone firewall policies, we can use a parallel parsing structure to automate the ingestion of industrial asset layouts. Understanding how firmware-level interaction rules require customized definitions of entry points and target execution goals directly informs our method for integrating the MITRE ATT&CK for ICS taxonomy. This structural alignment ensures that lateral credential reuse vectors and legacy protocol vulnerabilities can be automatically parsed into rigorous Datalog facts, bridging the gap between human-readable configuration data and automated risk assessments.

\begin{table}[htbp]
\centering
\caption{Structural Realignment of Threat Modeling Abstraction Layers}
\begin{tabular}{lll}
\toprule
\textbf{Analysis Variable} & \textbf{Micro Model (MIRAGE)} & \textbf{Macro Model (This Research)} \\
\midrule
Ingestion Target & Compressed Firmware Filesystems & JSON Asset Configuration Sheets \\
Parsing Mechanism & Binwalk Firmware Decomposition & Advanced LLM Parsing Framework \\
Data Propagation & Inter-Process Communication Flows & Zone-to-Zone Firewall Policies \\
Vulnerability Mapping & NIST NVD String Matches & Lateral Credential Reuse Tracking \\
Reasoning Rules & Supply Chain \& Execution Logic & Purdue Hierarchy \& OT Protocols \\
\bottomrule
\end{tabular}
\end{table}

\section{Architecture of High Level of Our Framework}
Within the context of our automated threat modeling pipeline, these optimization algorithms serve a crucial function at the boundary between logic processing and threat intelligence mapping. 

Once the automated parser synthesizes raw infrastructure configurations (\texttt{configurations.json}) into declarative Datalog clauses (\texttt{environmentNetwork.P}), the MulVAL reasoning engine applies its custom interaction rules (\texttt{attack1.P}) to build the foundational logical relationships. Rather than forcing a full, exhaustive enumeration of the resulting attack graph, which would introduce latency when scaling across expansive multi-subnet layouts, our architecture utilizes heuristic path discovery to quickly isolate the most critical propagation chains. 

By prioritizing the processing of high-probability paths (such as exploitless lateral movements enabled by shared administrative credentials), the optimization layer passes a refined threat matrix down to the enrichment phase. This ensures the pipeline can map vulnerabilities onto the MITRE ATT&CK for ICS framework rapidly and generate actionable hardening guidelines with minimal computational overhead.

\section{The Use of LLM to Generate Rules}

\subsection{State-of-the-Art Ingestion Paradigms: From Deterministic Parsing to Generative AI}
To resolve the structural disconnect between raw infrastructure states and logical verification engines, current research explores three primary translation paradigms:

\paragraph{Deterministic Parsers}
Deterministic parsers utilize custom, hardcoded scripts (written in languages like Python or Go) to parse structured configuration files and output text blocks formatted for Datalog engines. While these scripts guarantee absolute syntactic precision and predictable behavior, they lack semantic flexibility. If a vendor changes an inventory schema field from \texttt{software\_services} to \texttt{installed\_applications}, a deterministic script fails completely, requiring manual rewriting and maintenance.

\paragraph{Natural Language Processing (NLP) \& Tokenization Pipelines}
To introduce greater flexibility, earlier automated threat modeling frameworks implemented NLP techniques and heuristic token-matching algorithms. These pipelines extract keywords from unstructured documentation or software banners and map them to known system properties. While more resilient to minor vocabulary updates than rigid scripts, standard NLP models lack deep contextual awareness. They frequently struggle to maintain the strict structural formatting and exact variable alignments needed to compile executable Datalog scripts successfully.

\paragraph{Large Language Models (LLMs) for Logic Synthesis}
The latest paradigm shifts the translation task toward Large Language Models (LLMs) to automate formal logic synthesis. LLMs feature high contextual awareness and structural understanding, allowing them to process both semi-structured JSON blocks and unstructured network diagrams to generate precise Datalog clauses.
Instead of relying on rigid, hardcoded rules, an LLM-driven translation engine interprets data semantically. It can automatically reconcile naming variations, deduce implicit relationships (such as identifying that an asset operating as a \texttt{vdc01} node implies the presence of an Active Directory service), and map credential similarities across distinct hosts. This generative logic synthesis bridges the gap between raw, human-managed configuration logs and mathematically rigorous verification environments.

\subsection{Evaluating LLM Translation Accuracy and Validating Logical Soundness}
While LLMs introduce exceptional flexibility into the logic synthesis pipeline, utilizing generative AI within safety-critical industrial threat models introduces a critical challenge: ensuring absolute logical soundness and preventing syntax errors. Because language models operate probabilistically, they can introduce minor variations—such as formatting inconsistencies, variable mismatching, or hallucinated predicates—that prevent compile-time execution within strict Prolog engines.

To guarantee safety and validity before an automated script is sent to the reasoning engine, the pipeline must enforce a strict, multi-tiered validation matrix:
\begin{itemize}
    \item \textbf{Compile-Time Syntactic Check:} The output generated by the LLM is automatically processed through a linting layer to verify absolute compliance with XSB Prolog syntax rules. This includes confirming proper clause placement, verifying trailing periods, and validating that variables and atom casings align precisely with predefined schema parameters.
    \item \textbf{Deterministic Ground-Truth Mapping:} To eliminate the risk of logical hallucinations, synthesized predicates are evaluated against an immutable ground-truth schema matrix. This automated validation step verifies that generated predicates (such as \texttt{networkService/5} or \texttt{aclNW/4}) match authentic system parameters and network connections exactly.
\end{itemize}
By combining the semantic flexibility of LLM parsing with strict compile-time verification checks, the automated threat modeling pipeline achieves a reliable balance. It successfully automates the translation of real-world industrial infrastructure records into formal Datalog facts, delivering automated, mathematically sound risk assessments with zero operational disruption.


% ==========================================
% CHAPTER 5: IMPLEMENTATIONS
% ==========================================
\chapter{Implementations}
This chapter covers the technical deployment setup, including library definitions, API connection wrappers, the orchestration layout of the Python validation loop, and the rule parsing hooks written for the Prolog configuration files. Detailed structural components of the core modules will be systematically detailed as development steps conclude.


% ==========================================
% CHAPTER 6: TESTING: CASE STUDY
% ==========================================
\chapter{Testing: Case Study}
This chapter evaluates the practical execution of the pipeline over isolated test instances mimicking enterprise production environments. Computational performance metrics, model translation accuracies, syntax exception occurrences, and time optimizations are logged across scaling network layouts to verify model boundaries.


% ==========================================
% CHAPTER 7: CONCLUSIONS AND FUTURE WORKS
% ==========================================
\chapter{Conclusions and Future Works}
This final chapter evaluates the automated logic pipeline execution insights, assessing the real-world utility of semantic ingestion matrices alongside deterministic compilation layers. Future tracks investigate extending the static engine toward non-monotonic evaluation frameworks to dynamically trace changing transient network permissions over active cyber-physical feedback controls.


%% Fine dei capitoli normali, inizio dei capitoli-appendice (opzionali)
\appendix

\chapter{Custom Rule Sets Reference Blocks}
This appendix aggregates the core declarative Prolog predicates, interaction rules definitions, and template formatting specifications utilized across validation steps.

%% Parte conclusiva del documento; tipicamente per riassunto, bibliografia e/o indice analitico.
\backmatter

%% Bibliografia
\bibliographystyle{plain_\languagename}
\begin{thebibliography}{9}

\bibitem{Ou05} 
X. Ou, S. Govindavajhala, and A. W. Appel, 
"MulVAL: A Logic-based Network Security Analyzer," 
in \emph{Proceedings of the 14th USENIX Security Symposium}, Baltimore, MD, 2005.

\bibitem{TayouriSurvey} 
D. Tayouri, N. Baum, A. Shabtai, and R. Puzis, 
"A Survey of MulVAL Extensions and Their Attack Scenarios Coverage," 
\emph{arXiv preprint}, Department of Software and Information Systems Engineering, Ben-Gurion University of the Negev.

\bibitem{StanExtended} 
O. Stan, R. Bitton, M. Ezrets, M. Dadon, Y. Elovici, A. Shabtai, M. Inokuchi, Y. Ohta, Y. Yamada, and T. Yagyu, 
"Extending Attack Graphs to Represent Cyber-Attacks in Communication Protocols and Modern IT Networks," 
\emph{Preprint}, Ben-Gurion University of the Negev / NEC Corporation.

\bibitem{SunderOT} 
G. Sunder, A. S. Colletto, S. Raimondi, C. Basile, A. Viticchié, and A. Aliberti, 
"Enhancing OT Threat Modelling: An Effective Rule-Based Approach for Attack Graph Generation," 
Politecnico di Torino / AlphaWaves S.r.l., Turin, Italy.

\bibitem{KonstaSurvey} 
A.-M. Konsta, A. L. Lafuente, B. Spiga, and N. Dragoni, 
"Survey: Automatic generation of attack trees and attack graphs," 
\emph{Academic Press / Journal of Systems and Software}.

\bibitem{AbdellatifAAPDEM} 
T. M. Abdellatif, A. Hamed, S. R. Majeed, and A. A. Seyam, 
"AAPDEM: An Intelligent Model for Automated Attack Path Discovery and Exploitation in Modern Networks," 
\emph{IEEE Access}, vol. 13, 2025.

\bibitem{StammetGNN} 
C. Stammet, P. Dotti, U. Ultes-Nitsche, and A. Fischer, 
"Analyzing Büchi Automata with Graph Neural Networks," 
University of Fribourg / University of Applied Sciences and Arts Western Switzerland.

\bibitem{TayouriMIRAGE} 
D. Tayouri, T. Nachum, and A. Shabtai, 
"MIRAGE: Multi-Binary Image Risk Assessment With Attack Graph Employment," 
in \emph{IEEE Transactions on Dependable and Secure Computing}, vol. 23, no. 2, pp. 2336-2353, March-April 2026.

\end{thebibliography}

\end{document}
